Beyond strategic design, the day-to-day performance of a \ac{bss} depends critically on the ability to anticipate and manage operational imbalances. As detailed in Chapters~\ref{ch:mobility_and_battery_modeling} and ~\ref{ch:failure_forecasting}, two challenges have become increasingly important. The first is to predict where bikes will move across the network and how their battery levels will evolve, since both factors strongly affect rebalancing and bike availability. The second is to improve the efficiency of maintenance operations by identifying when key components are likely to fail, so that repairs can be planned in advance in order to use resources more effectively.

Part~\ref{pt:predictive_operations} develops two parallel predictive frameworks using operational data from Barcelona’s \ac{bss}, \textit{Bicing}. After introducing \textit{Bicing} in Chapter~\ref{ch:use_case_bicing}, Chapter~\ref{ch:results_mobility_and_battery_modeling} presents the first framework, which combines a Markov chain mobility model with a physics-informed battery dynamics function to forecast bike locations and battery levels. Then, Chapter~\ref{ch:results_failure_forecasting} presents the second framework, which uses \ac{sa} and \ac{ml} techniques to estimate failure risks for key components of both \acp{m-b} and \acp{e-b}. The material in this part is based on two published articles: Bassolas et al.~\cite{Bassolas2025} and Grau-Escolano et al.~\cite{Grau-Escolano2024}, published in \textit{Scientific Reports} and \textit{EPJ Data Science} respectively.